\section{Danh sách, bảng và hình}
\label{sec:khoi}

\subsection{Danh sách}
\begin{ma}
\begin{itemize}        % gạch đầu dòng; dùng enumerate để đánh số
  \item Ý thứ nhất
  \item Ý thứ hai
\end{itemize}
\end{ma}

\begin{meo}
Khi đặt con trỏ tại một dòng \lenh{item}, bạn có thể nhấn \phim{Alt}+\phim{↑}
hoặc \phim{Alt}+\phim{↓} để đổi chỗ cả mục đó (bao gồm các dòng con bên trong)
với mục liền kề.
\end{meo}

\subsection{Bảng}
Cách tạo bảng nhanh nhất là gõ \texttt{\textbackslash table} rồi nhấn
\phim{Enter} để chèn một bảng mẫu, sau đó chỉnh sửa trực quan bằng hộp thoại
bảng (qua biểu tượng ✎ trong Mục lục):

\anh[0.62\linewidth]{popup-bang}{Ô sửa bảng: nhập nội dung trực tiếp, kéo ⠿ để
đổi vị trí hàng hoặc cột, bấm \nut{+ Row} để thêm hàng}

\begin{itemize}
  \item Kéo biểu tượng ⠿ ở đầu hàng hoặc đầu cột để sắp xếp lại thứ tự; hệ
    thống sẽ hiển thị trước vị trí mới trong lúc bạn kéo.
  \item Nhấp vào biểu tượng căn lề trên đầu cột để chuyển đổi giữa căn trái,
    giữa và phải.
  \item Điền \nut{Caption} (chú thích) và \nut{Label} (nhãn) để thuận tiện cho
    việc tham chiếu tới bảng trong bài viết.
  \item Trong mục \nut{More}, chọn \nut{Cleaner rules} để định dạng đường kẻ
    bảng gọn gàng theo chuẩn sách in.
\end{itemize}

\begin{meo}
Ngay trong trình soạn thảo mã, khi đặt con trỏ vào một ô bất kỳ của bảng, bạn
có thể dùng \phim{Alt}+\phim{←} hoặc \phim{Alt}+\phim{→} để di chuyển cả cột đó
sang trái hoặc sang phải.
\end{meo}

\subsection{Hình ảnh}
Bạn có thể \textbf{dán trực tiếp} bằng cách chụp màn hình rồi nhấn
\phim{Ctrl}+\phim{V} vào trình soạn thảo, hoặc kéo thả tệp ảnh từ máy tính vào.
SumiTeX sẽ tự động xử lý các bước sau:
\begin{enumerate}
  \item Tải ảnh lên thư mục \texttt{images/} với tên tệp hợp lệ.
  \item Chèn sẵn một khối hình hoàn chỉnh và bôi đen phần chú thích để bạn gõ
    đè ngay.
  \item Bổ sung gói \texttt{graphicx} nếu tài liệu chưa khai báo.
\end{enumerate}
Ngoài ra, một thanh công cụ nhỏ phía trên khối hình sẽ giúp bạn chọn nhanh độ
rộng theo các mức ¼ · ½ · ¾ · Full.

\begin{chuy}
Vì \LaTeX{} không hỗ trợ trực tiếp các định dạng WebP, GIF và SVG, SumiTeX sẽ
tự động chuyển đổi chúng sang định dạng PNG khi bạn dán vào tài liệu.
\end{chuy}

\begin{thuchanh}[Bài 2, 3, 4]
Hãy hoàn thành lần lượt các tệp \texttt{bai-02.tex} (tạo danh sách),
\texttt{bai-03.tex} (tạo bảng) và \texttt{bai-04.tex} (dán một ảnh bất kỳ,
chỉnh độ rộng ½ và thêm chú thích).
\end{thuchanh}